% --- Fate's Edge SRD --- Section 11: Backlash, Obligation, and Ritual Consequences ---
\section{Backlash, Obligation, and Ritual Consequences}
\label{sec:backlash-obligation}

Magic is never free. This section codifies the three main risk currencies: \textbf{Backlash} (for free casting), \textbf{Obligation} (for Rites and Patron dealings), and \textbf{Ritual Consequences} (for invokers and universal rituals).

\subsection{Backlash (Free Casting)}
Backlash is the unpredictable fallout from shaping raw magic. It occurs when a casting roll (Weave or Cast) yields a \textbf{Partial} or \textbf{Miss}, or on a \textbf{Hit that shows two or more 1s} (minor backlash alongside success).

\paragraph{Backlash Menu (choose one on Partial, two on Miss)}
\begin{itemize}
  \item \textbf{Position Shift:} Your Position worsens by one step (Dominant $\to$ Controlled $\to$ Desperate) for your current or next action.
  \item \textbf{Fleeting Harm/Condition:} Minor injury (1 Fatigue or Harm 1) or a temporary Condition (Dazed, Singed, Disoriented).
  \item \textbf{Exposure/Noise:} Draws unwanted attention – guards alerted, a door slams, a rival notices.
  \item \textbf{Resource Drain:} Lose a component, a tool wears out, or you mark 1 Fatigue.
  \item \textbf{Collateral Spark:} A nearby ally, object, or environmental feature is threatened (GM describes).
\end{itemize}
All backlash is \textbf{colored by the dominant element} (see Elemental Backlash table, §10). The GM spends Story Beats (generated by 1s) to apply these consequences.

\subsection{Obligation (Rites and Patrons)}
Obligation measures a Patron's claim on a character. It is tracked per Patron as a clock (default 6 segments). 

\paragraph{Marking Obligation}
\begin{itemize}
  \item \textbf{Invoking a Rite (Runekeeper):} Mark the Rite's listed Obligation (typically +1).
  \item \textbf{Pushing a Rite:} Mark an additional +1 Obligation.
  \item \textbf{Cross-Patron Interference:} Using Rites from two different Patrons in the same scene marks +1 Obligation to the second Patron.
\end{itemize}

\paragraph{Obligation Capacity and Overflow}
Obligation Capacity = Spirit + Presence. If total Obligation exceeds Capacity:
\begin{itemize}
  \item Mark 1 Fatigue per segment over Capacity.
  \item You cannot invoke new Rites or perform Patron rituals until Obligation is reduced below Capacity.
\end{itemize}

\paragraph{Clearing Obligation}
\begin{itemize}
  \item \textbf{Off‑screen:} Perform an Act of Service aligned with the Patron (provide an Intricate Description). Clear 1 segment.
  \item \textbf{On‑screen (Focused Devotion talent, 3–5 XP):} Once per scene, spend 1 Boon to clear 1 segment (cannot clear segments accrued in the same scene).
\end{itemize}

\subsection{Ritual Consequences (Invokers and Universal Rituals)}
Rituals (slow, significant‑time workings) carry their own risks beyond ordinary casting.

\paragraph{Invoker Rituals (via Symbol)}
\begin{itemize}
  \item \textbf{Ritual Completion:} Always marks \textbf{+1 Obligation} (attention cost).
  \item \textbf{Crack the Seal (Instant Cast):} Mark \textbf{+2 Obligation} (+3 for High Rites); the Symbol becomes \textsc{Compromised} (inert until restored). The GM may immediately spend 1 SB on‑theme.
  \item \textbf{Restoring a Symbol:} Downtime action + test (DV 3, Arcana or Lore). Success returns to Maintained; shaky success leaves it Neglected (future rituals +1 Obligation).
\end{itemize}

\paragraph{Universal Rituals (any character)}
\begin{itemize}
  \item \textbf{Failure / Partial:} The GM may impose a consequence from the Backlash Menu (see above) or a narrative twist (wrong spirit answers, ward backfires).
  \item \textbf{Components:} Rare components may be consumed on a Partial or Miss.
\end{itemize}

\paragraph{Symbol Overload and Rivalries}
\begin{itemize}
  \item \textbf{Multi‑Symbol Overload:} Carrying 4+ different Symbols causes the first Invoker ritual each scene to mark +1 additional Obligation.
  \item \textbf{Rival Symbols:} Displaying a rival Patron's Symbol while invoking another worsens Position by 1 and marks +1 Obligation to the current Patron. On any rolled 1, the GM may trigger a Patron omen/glitch.
\end{itemize}

\subsection{Interruption and Counterplay}
\begin{itemize}
  \item \texttt{[COUNTER]} can interrupt a Weave, Cast, Rite, or ritual within its casting window (DV by fiction).
  \item \texttt{[DISPEL]} ends or suppresses ongoing magical effects (DV by fiction).
  \item \texttt{[UNWARD]} suppresses/dismisses \texttt{[WARD]}s; \texttt{[BANISH]} and \texttt{[WARD]} interact with Outsiders per §10.
\end{itemize}

\subsection{GM Budget Dials}
\begin{itemize}
  \item \textbf{SB Budgets:} Use scene/session limits (see Complications, §7).
  \item \textbf{Obligation Pace:} Expect 2–4 segments marked per session for an active Runekeeper.
  \item \textbf{Backlash Visibility:} Favor big, legible consequences over many minor pinpricks.
\end{itemize}

\subsection{Micro‑Examples}

\paragraph{Free Casting Backlash}
\emph{Lyra casts a fire bolt (Weave + Cast). Partial on the Cast roll with two 1s. The GM spends SB: "The bolt hits, but the flash blinds you – your next action is Desperate. Also, the smoke triggers a hanging lantern; the room starts catching fire (Clock: Fire Spread [4])."}

\paragraph{Runekeeper Push and Debt}
\emph{Kael invokes the Sealed Gate's \textit{Circle of Denial} \texttt{[WARD]} and Pushes It to harden the ring. He marks +1 Obligation for the Rite and +1 for the push (total +2). Later, a demon tests the ring – the GM uses the \texttt{[WARD]} rules (DV = demon's Cap). On a Hit, the demon crosses, adding +DV to its Leash.}

\paragraph{Crack the Seal Under Fire}
\emph{Vex presents Ikasha's Symbol and Cracks the Seal to lay an instant shadow lane. Symbol becomes Compromised; she marks +2 Obligation. The GM immediately spends 1 SB to dim all lights (panic). Later, during downtime, Vex restores the Symbol with an Arcana test (DV 3) – a shaky success leaves it Neglected until she performs a full rite of cleaning.}